\section {The spy (Information theory)}

A spy in an enemy country wants to send information home using the enemy's daily morning radio transmission, which is 15 bits long.
Every morning the transmission is random. The spy can get into the radio station a few minutes before the transmission, look at the 15 bits that are about to be
sent, and either leave them as they are or flip a single bit (flipping more would make the transmission look corrupt).
His operators listen to the transmission and decode it according to a method agreed on in advance.
How much information can the spy send every day?

\subsection{Solution}

The spy can send 4 bits, that is, one of 16 messages, every day.

He cannot do better: every morning he has exactly 16 choices (flip one of the 15 bits, or flip nothing),
and the operators see only the result, so at most 16 different messages can be told apart.

Here is how to reach 16. Number the positions of the bits $1$ to $15$ and think of each position number as $4$ bits.
The operators decode the transmission by computing the exclusive or (xor) of the positions of all the bits that are $1$.
That is a number between $0$ and $15$, and it is the message.

The spy computes the same value $v$ for the transmission he finds. He wants to send the message $m$.
Flipping the bit at position $k$ adds $k$ to the set of positions of the $1$ bits or removes it, and either way the value becomes $v \oplus k$.
So if $v = m$ he leaves the transmission alone, and otherwise he flips the bit at position $k = v \oplus m$, which is between $1$ and $15$.
The value becomes $v \oplus v \oplus m = m$.

This is the Hamming code of length 15 in disguise. Compare with the riddle of the coins on the chess board, which uses the same trick.
